\documentclass[11pt,a4paper]{article}
\usepackage[utf8]{inputenc}
\usepackage{geometry}
\usepackage{xcolor}
\usepackage{enumitem}
\usepackage{hyperref}
\usepackage{fontawesome5}
\usepackage{titlesec}
\usepackage{pifont}

% Disable hyphenation
\usepackage{hyphenat}
% Page layout
\geometry{margin=1in}
\setlength{\parindent}{0pt}
\setlength{\parskip}{4pt}

% Color scheme
\definecolor{text}{RGB}{50, 50, 50}

% Section formatting
\titleformat{\section}
{\large\bfseries}
{}
{0em}
{}[\titlerule]

% Hyperlink settings
\hypersetup{
    hidelinks                   % 隐藏所有链接边框和背景色
}

% Custom commands
\newcommand{\cvitem}[2]{\textbf{#1} \hfill \textit{#2}}
\newcommand{\skillitem}[2]{\textbf{#1:} #2}
\newcommand{\contact}[2]{\faIcon{#1} \hspace{0.5em} #2}

\begin{document}

% Header
\begin{center}
    {\Huge\bfseries Xiaobin Ren}\\[4pt]
    % Contact information
    \contact{envelope}{\href{mailto:1155239443@link.cuhk.edu.hk}{\texttt{1155239443@link.cuhk.edu.hk}}} \quad
    \contact{phone}{+852 92870301} \quad
    \contact{map-marker-alt}{Hong Kong, China} \\
    \contact{globe}{{\href{https://georgexiaobinren.github.io/}{\texttt{georgexiaobinren.github.io}}}} \quad
    \contact{github}{\href{https://github.com/GeorgeXiaobinRen}{\texttt{github.com/GeorgeXiaobinRen}}}
\end{center}

% Education
\section{Education}
\cvitem{M.S. in Mathematics}{Aug. 2025 - Jun. 2026(expected)}\\
\ding{226} \textbf{Chinese University of Hong Kong}\\
Relevant Courses: Stochastic Analysis, Mathematical Image Processing, Machine Learning, Artificial Intelligence\\

\cvitem{B.S. in Mathematics}{Sep. 2021 - Jun. 2025}\\
\ding{226} \textbf{Hunan University} (Project 985)\\
GPA of compulsory courses: 3.46/4.0 (85.44/100)\\
Relevant Courses: Numerical Solution of Differential Equations, Optimization Theory, Numerical Analysis

% Technical Skills
\section{Technical Skills}
\begin{itemize}[leftmargin=*,nosep]
    \item \skillitem{Programming Languages}{Python, MATLAB, R, C++, MySQL}
    \item \skillitem{Typesetting Tools}{\LaTeX, Markdown, Typst, Word}
    \item \skillitem{Mathematical Software}{Mathematica, Maple}
    \item \skillitem{Data Analysis}{NumPy, Pandas, SciPy}
    \item \skillitem{Machine Learning}{PyTorch, Scikit-learn}
    \item \skillitem{Tools}{Git, Jupyter Notebook, Linux, Docker}
\end{itemize}

% Projects
\section{Research Projects}
\cvitem{Application of Deep Learning Method in Numerical Computation}{\textit{2023 - 2025}}
\begin{itemize}[leftmargin=*,nosep]
    \item China University Student Innovation Training Program (SIT) (National-level)
    \item Position: Team member
    \item Instituter: Prof. Wenfan Yi (Hunan University)
    \item Responsibilities:
    \begin{itemize}[leftmargin=*,nosep]
        \item Learned deep learning methods including ResNet, PINNs, SINDy, etc.
        \item Conducted in-depth research on the theoretical foundations of the KAN network and analyzed its potential advantages.
        \item Designed the architecture and configure parameters, and performed comparative experiments to evaluate the performance of KAN with other networks on
different problems.
        \item Participated in completing academic paper
    \end{itemize}
\end{itemize}
\vspace{0.5cm}
\cvitem{Meshless Computational Methods for Financial Derivatives Problems}{\textit{2025 - 2026}}
\begin{itemize}[leftmargin=*,nosep]
    \item Independent study project
    \item Position: Team member
    \item Instituter: Prof. Benny Y. C. Hon (Chinese University of Hong Kong)
    \item Responsibilities:
    \begin{itemize}[leftmargin=*,nosep]
        \item Studied classical numerical methods for fractional differential equations (FPDEs) and Volterra integral equations
        \item Converted fractional differential operators and integrals into approximate formats that can be handled by automatic differentiation
        \item Considered applying Monte Carlo sampling to different problems (especially the high-dimensional problems and problems defined on non-rectangular domains)
    \end{itemize}
\end{itemize}

% Research Experience
%\section{Research Experience}
%\cvitem{}
%\textit{September 2023 - Present}
%\begin{itemize}[leftmargin=*,nosep]
%    \item Assisted professor in computational mathematics research projects
%    \item Developed algorithms for numerical linear algebra applications
%    \item Conducted literature reviews and contributed to research papers
%\end{itemize}

% Work Experience
\section{Work Experience}
\cvitem{Data Analysis Intern}{Jul. 2023 - Aug. 2023}
\begin{itemize}[leftmargin=*,nosep]
    \item Tianjin Xinghu Technology Development Co., LTD
    \item Responsibilities:
    \begin{itemize}[leftmargin=*,nosep]
        \item Extracted relevant data from various sources, including company databases, online surveys, and market research, and conducted forecasts on key metrics such as trends in product price changes.
        \item Utilized tools like Excel and Python to analyze collected data, identify trends and patterns, and generate analytical reports to support decision-making.
    \end{itemize}
\end{itemize}

% Awards \& Honors
\section{Awards \& Honors}
\begin{itemize}[leftmargin=*,nosep]
    \item Second-Class Scholarship in the 2021-2022 Academic Year (University-Level), 2022
    \item Third Prize in the China Undergraduate Mathematical Contest in Modeling (Provincial-Level), 2023

\end{itemize}

% Languages
\section{Languages}
\begin{itemize}[leftmargin=*,nosep]
    \item English (fluent, IELTS 7.0, CET6 580)
    \item Mandarin (native)
    \item Dialects: Cantonese (a little), Wu dialect (native)
\end{itemize}

% References
%\section{Publications}
%\begin{itemize}[leftmargin=*,nosep]
%    \item Author, X., \textbf{Xiaobin Ren}, et al. ``Title of Paper.'' \textit{Journal Name}, Year.
%    \item \textbf{Xiaobin Ren}, Co-author, Y. ``Another Paper Title.'' \textit{Conference Proceedings}, Year.
%\end{itemize}


\end{document}